% ==========================================================================
%  Chapter 55 — Seismic Analysis Infrastructure
% ==========================================================================

\chapter{Seismic Analysis Infrastructure}
\label{ch:seismic-analysis-infra}

This chapter documents the comprehensive seismic analysis stack added
to the \texttt{fall\_n} library.  The infrastructure includes:
%
\begin{enumerate}
  \item a \textbf{GroundMotionRecord} parser for earthquake acceleration
        files,
  \item a \textbf{DamageCriterion} framework for identifying the most
        damaged elements (designed as a strategy/injection point),
  \item a \textbf{FiberHysteresisRecorder} observer for tracking
        stress--strain histories in extreme fibers,
  \item the \texttt{section\_snapshots()} virtual in
        \texttt{StructuralElement} --- the enabling infrastructure
        for generic fiber access,
  \item a complete \textbf{five-story RC building example} under the
        El~Centro 1940 record with corotational elements, and
  \item a full test suite validating the infrastructure.
\end{enumerate}

All code targets C++23 and follows the \texttt{fall\_n} conventions:
external polymorphism for type erasure, strategy-pattern injection,
observer-based I/O, and PETSc for the linear algebra backend.


% ══════════════════════════════════════════════════════════════════════════
\section{Earthquake Ground Motion Input}
\label{sec:seismic-ground-motion}

% ── 1.1 ──────────────────────────────────────────────────────────────────
\subsection{Data Format and the El Centro Record}

The library ships a reference earthquake record in
\texttt{data/input/earthquakes/el\_centro\_1940\_ns.dat},
stored in two-column whitespace-delimited format:
%
\begin{lstlisting}[language={},basicstyle=\ttfamily\small]
# El Centro 1940 NS -- Imperial Valley earthquake
# Source: PEER NGA-West2 (RSN6)
# 406 points, dt = 0.02 s, PGA = 0.3194 g at t = 4.80 s
# Units: time [s], acceleration [g]
0.0000    0.00630
0.0200    0.00364
0.0400    0.00099
...
\end{lstlisting}
%
Lines beginning with \texttt{\#} are comments.  Each data line
contains a time value and the corresponding ground acceleration.
The El~Centro 1940 North--South component has 406~points
spanning $t \in [0, 10]$~s with $\Delta t = 0.02$~s and a peak
ground acceleration (PGA) of approximately $0.3194g$ at
$t \approx 4.80$~s.


% ── 1.2 ──────────────────────────────────────────────────────────────────
\subsection{The \texttt{GroundMotionRecord} Class}

\texttt{src/model/GroundMotionRecord.hh} provides a lightweight parser
and accessor class in the \texttt{fall\_n} namespace.

\paragraph{Factory methods.}
Three static constructors handle different input scenarios:
%
\begin{itemize}
  \item \texttt{from\_file(path, scale)} --- parses a two-column file,
        multiplying all acceleration values by \texttt{scale}.
  \item \texttt{from\_single\_column(path, dt, scale)} --- for files
        containing only acceleration values at a fixed $\Delta t$.
  \item \texttt{from\_vectors(times, accels, name)} --- in-memory
        construction from explicit \texttt{std::vector<double>}.
\end{itemize}

\paragraph{Accessors.}
\begin{center}
\begin{tabular}{ll}
  \toprule
  Method & Description \\
  \midrule
  \texttt{pga()} & Peak ground acceleration ($\max |a_g|$) \\
  \texttt{time\_of\_pga()} & Time instant of PGA \\
  \texttt{dt()} & Average time step \\
  \texttt{duration()} & Total record duration ($t_{\max}$) \\
  \texttt{num\_points()} & Number of data points \\
  \texttt{name()} & Descriptive record name \\
  \bottomrule
\end{tabular}
\end{center}

\paragraph{TimeFunction conversion.}
The method \texttt{as\_time\_function()} returns a callable
\texttt{TimeFunction} that performs piecewise-linear interpolation
over the data points, with zero-padding beyond the record boundaries.
This is the key bridge into the dynamic analysis pipeline:
%
\begin{lstlisting}[language=C++,basicstyle=\ttfamily\small]
auto record = fall_n::GroundMotionRecord::from_file(
    "el_centro_1940_ns.dat", 9.81);   // scale g -> m/s^2
auto ag = record.as_time_function();   // TimeFunction
\end{lstlisting}

\paragraph{Convenience factory.}
The free function \texttt{make\_ground\_motion\_from\_file(path, dir,
scale)} combines parsing with \texttt{GroundMotionBC} construction
in a single call:
%
\begin{lstlisting}[language=C++,basicstyle=\ttfamily\small]
auto gm = fall_n::make_ground_motion_from_file(
    eq_path, /*direction=*/0, /*scale=*/9.81);
bcs.add_ground_motion(std::move(gm));
\end{lstlisting}
%
The direction index $d \in \{0,1,2\}$ corresponds to
$\{X, Y, Z\}$.  The resulting \texttt{GroundMotionBC} is consumed
by \texttt{BoundaryConditionSet::add\_ground\_motion()}, which the
\texttt{DynamicAnalysis} driver queries to construct the
pseudo-force vector $\mathbf{M}\hat{\mathbf{g}}\,a_g(t)$.


% ══════════════════════════════════════════════════════════════════════════
\section{The \texttt{section\_snapshots()} Virtual}
\label{sec:seismic-section-snapshots}

A recurring limitation in the library was the inability to query
fiber-level state \emph{generically} through the
\texttt{StructuralElement} type-erased wrapper.  Observer code that
needed fiber strains had to downcast via \texttt{as<T>()} to each
concrete beam type --- an unmaintainable and fragile approach.

\subsection{Design Decision}

We added a single virtual method to the
\texttt{StructuralElement::Concept} base:
%
\begin{lstlisting}[language=C++,basicstyle=\ttfamily\small]
virtual std::vector<SectionConstitutiveSnapshot>
    section_snapshots() const { return {}; }
\end{lstlisting}
%
The \texttt{Model<T>} override uses a \texttt{if constexpr
(requires)} check to dispatch statically:
%
\begin{lstlisting}[language=C++,basicstyle=\ttfamily\small]
std::vector<SectionConstitutiveSnapshot>
section_snapshots() const override {
    if constexpr (requires { element_.sections(); }) {
        const auto& secs = element_.sections();
        std::vector<SectionConstitutiveSnapshot> result;
        result.reserve(secs.size());
        for (const auto& s : secs)
            result.push_back(s.section_snapshot());
        return result;
    } else {
        return {};
    }
}
\end{lstlisting}
%
Elements without sections (e.g., continuum elements) return an
empty vector.  Beam elements return one snapshot per Gauss point,
each containing beam properties and (optionally) fiber-level data.
Shell elements return shell section data.

\subsection{SectionConstitutiveSnapshot Anatomy}

\begin{lstlisting}[language=C++,basicstyle=\ttfamily\small]
struct SectionConstitutiveSnapshot {
    std::optional<BeamSectionConstitutiveSnapshot>  beam;
    std::optional<ShellSectionConstitutiveSnapshot> shell;
    std::span<const FiberSectionSample>             fibers;

    bool has_beam()   const noexcept;
    bool has_shell()  const noexcept;
    bool has_fibers() const noexcept;
};
\end{lstlisting}
%
Each \texttt{FiberSectionSample} contains:
\texttt{y}, \texttt{z} (position), \texttt{area},
\texttt{strain\_xx}, and \texttt{stress\_xx}.
The \texttt{fibers} span borrows memory from the fiber section's
internal storage (zero-copy).

This single virtual eliminates all downcast-based fiber access in:
\begin{itemize}
  \item \texttt{FiberRecorder::on\_step()} in \texttt{Observers.hh},
  \item \texttt{ElementRecorder::on\_step()} in \texttt{Observers.hh},
  \item the new \texttt{DamageCriterion} and
        \texttt{FiberHysteresisRecorder} modules.
\end{itemize}


% ══════════════════════════════════════════════════════════════════════════
\section{Damage Criterion Framework}
\label{sec:seismic-damage-criterion}

The identification of the most damaged elements in a structure
requires a pluggable damage metric.  We implement this as a
classic \emph{strategy pattern} with a virtual base class,
providing an injection point for user-defined criteria.

\subsection{Class Hierarchy}

\begin{center}
\begin{tabular}{lp{8cm}}
  \toprule
  Class & Role \\
  \midrule
  \texttt{DamageCriterion} & Abstract base.
    Pure virtuals: \texttt{name()}, \texttt{evaluate\_element()},
    \texttt{evaluate\_fibers()}, \texttt{clone()} \\
  \texttt{MaxStrainDamageCriterion} & Default implementation.
    Damage index $D = \max_i |\varepsilon_i| / \varepsilon_{\text{ref}}$
    over all fibers \\
  \texttt{CallableDamageCriterion} & Wraps a \texttt{std::function}
    for quick user injection \\
  \bottomrule
\end{tabular}
\end{center}

\subsection{Evaluation Protocol}

\texttt{evaluate\_element(elem, idx, u)} returns an
\texttt{ElementDamageInfo} containing:
%
\begin{itemize}
  \item \texttt{element\_index} --- global index in the model,
  \item \texttt{damage\_index} --- scalar damage measure $D \in [0, \infty)$,
  \item \texttt{max\_gp} --- Gauss point with highest fiber damage,
  \item \texttt{max\_fiber} --- fiber index within that GP.
\end{itemize}

\texttt{evaluate\_fibers(elem, idx, u)} returns a
\texttt{std::vector<FiberDamageInfo>}, one per fiber, with per-fiber
position $(y,z)$, strain $\varepsilon$, stress $\sigma$, and
individual damage index.

Both methods receive the raw PETSc \texttt{Vec u}
(current displacement) for criteria that need element-local
state beyond what \texttt{section\_snapshots()} provides.

\subsection{MaxStrainDamageCriterion}

The default criterion evaluates
%
\begin{equation}
  D_i = \frac{|\varepsilon_{xx,i}|}{\varepsilon_{\text{ref}}},
  \qquad
  D = \max_i D_i
\end{equation}
%
where $\varepsilon_{\text{ref}}$ is a user-supplied reference strain
(typically the yield strain $\varepsilon_y = f_y / E_s$).
Values $D > 1$ indicate yielding.

\subsection{CallableDamageCriterion and \texttt{make\_damage\_criterion}}

For rapid prototyping, users can inject a lambda:
%
\begin{lstlisting}[language=C++,basicstyle=\ttfamily\small]
auto custom = fall_n::make_damage_criterion(
    "ParkAng",
    [](const StructuralElement& elem,
       std::size_t idx, Vec u) -> ElementDamageInfo {
        // user logic here
        return {idx, damage_index};
    });
\end{lstlisting}
%
The \texttt{CallableDamageCriterion} owns the function object and
supports \texttt{clone()} for deep copying.


\subsection{DamageTracker Observer}

\texttt{DamageTracker<ModelT>} is a composable observer that
evaluates the damage criterion at configurable intervals:
%
\begin{itemize}
  \item \texttt{on\_step()}: evaluates all elements, maintains a
        sorted ranking and a peak (envelope) ranking.
  \item \texttt{on\_analysis\_end()}: prints the top-$N$ most
        damaged elements across the entire analysis.
\end{itemize}
%
\begin{lstlisting}[language=C++,basicstyle=\ttfamily\small]
fall_n::MaxStrainDamageCriterion criterion(eps_yield);
fall_n::DamageTracker<ModelT> tracker(criterion, /*interval=*/10);
\end{lstlisting}


% ══════════════════════════════════════════════════════════════════════════
\section{Fiber Hysteresis Recorder}
\label{sec:seismic-fiber-hysteresis}

The \texttt{FiberHysteresisRecorder<ModelT>} observer discovers
fibers at the first step, classifies them by material type
(concrete vs.\ steel), and records complete stress--strain
time-histories $(\varepsilon_{xx}(t), \sigma_{xx}(t))$ for later
hysteresis plotting.

\subsection{Architecture}

\paragraph{Fiber classification.}
A user-injected \texttt{FiberClassifier} function maps each
\texttt{FiberSectionSample} to a \texttt{FiberMaterialClass}:
%
\begin{lstlisting}[language=C++,basicstyle=\ttfamily\small]
enum class FiberMaterialClass { Concrete, Steel, Unknown };

using FiberClassifier = std::function<
    FiberMaterialClass(const FiberSectionSample&,
                       std::size_t elem_idx)>;
\end{lstlisting}
%
The default classifier uses a heuristic based on the initial
tangent stiffness (stress/strain ratio at the snapshot), but
users can override it.

\paragraph{Discovery phase.}
On the first time step, the observer iterates all target elements,
calls \texttt{section\_snapshots()}, and creates a
\texttt{FiberTrack} for each fiber with classification metadata.

\paragraph{Recording.}
Each subsequent step appends $(\varepsilon, \sigma)$ to the
corresponding track.

\paragraph{Post-analysis.}
\texttt{write\_hysteresis\_csv(basename)} exports two CSV files:
\texttt{\_concrete.csv} and \texttt{\_steel.csv}, each with columns:
\texttt{elem, gp, fiber, y, z, peak\_damage, eps\_1, sig\_1,
eps\_2, sig\_2, \ldots}

\subsection{Integration with DamageCriterion}

After the analysis, the recorder internally runs the configured
\texttt{DamageCriterion} (default: \texttt{MaxStrainDamageCriterion})
to sort fibers by peak damage.  The \texttt{top\_fibers(mat\_class)}
method returns the top-$N$ most damaged tracks of a given material
class, providing exactly the ``extreme fibers'' requested for
hysteresis plots.


% ══════════════════════════════════════════════════════════════════════════
\section{Observer Wiring: FiberRecorder and ElementRecorder}
\label{sec:seismic-observer-wiring}

Prior to this work, \texttt{FiberRecorder::on\_step()} and
\texttt{ElementRecorder::on\_step()} in \texttt{Observers.hh}
contained \texttt{TODO} stubs.  Both now use
\texttt{section\_snapshots()} to populate their internal
data structures:

\paragraph{FiberRecorder.}
Each step, the recorder iterates elements, obtains snapshots with
\texttt{has\_fibers()}, and stores per-fiber strain/stress from
the \texttt{FiberSectionSample} span.

\paragraph{ElementRecorder.}
Each step, the recorder extracts beam section properties
($E$, $G$, $A$, $I_y$, $I_z$, $J$) from
\texttt{BeamSectionConstitutiveSnapshot} or shell properties
from the \texttt{ShellSectionConstitutiveSnapshot}.


% ══════════════════════════════════════════════════════════════════════════
\section{Five-Story RC Building Under El Centro}
\label{sec:seismic-building-example}

\texttt{main\_seismic.cpp} implements a complete seismic analysis
example:

\subsection{Structural Configuration}

\begin{center}
\begin{tabular}{ll}
  \toprule
  Parameter & Value \\
  \midrule
  Stories & 5 \\
  Story height & 3.0~m \\
  Bay width & 5.0~m \\
  Column section & $0.50 \times 0.50$~m RC \\
  Beam section & $0.30 \times 0.50$~m RC \\
  Slab thickness & 0.15~m \\
  Concrete $f'_c$ & 28~MPa (Kent--Park) \\
  Steel $f_y$ & 420~MPa (Menegotto--Pinto) \\
  \bottomrule
\end{tabular}
\end{center}

\subsection{Element Formulations}

\begin{itemize}
  \item \textbf{Columns and beams:}
    \texttt{BeamElement<TimoshenkoBeam3D, 3, beam::Corotational>}
    with fiber sections built via \texttt{make\_rc\_column\_section}
    and \texttt{make\_rc\_beam\_section}.
  \item \textbf{Slabs:}
    \texttt{CorotationalMITC4Shell<>} (MITC4 assumed-strain shell
    with corotational kinematics) using
    \texttt{MindlinShellMaterial}.
\end{itemize}

The combination of corotational beams with MITC4 shells avoids
the membrane locking of standard displacement-based shells and
maintains numerical stability under large rotations.

\subsection{Loading Sequence}

\begin{enumerate}
  \item \textbf{Gravity ramp}: self-weight applied as a quasi-static
    ramp over the first time steps.
  \item \textbf{Ground motion}: the El~Centro 1940 NS record
    (scaled to m/s$^2$) applied in the $X$-direction via
    \texttt{BoundaryConditionSet::add\_ground\_motion()}.
    The \texttt{DynamicAnalysis} driver pre-computes the
    influence vector $\mathbf{M}\hat{\mathbf{g}}$ and applies
    the pseudo-force $a_g(t)\,\mathbf{M}\hat{\mathbf{g}}$
    at each time step.
\end{enumerate}

\subsection{Observer Pipeline}

The example configures a six-component
\texttt{CompositeObserver}:
%
\begin{enumerate}
  \item \texttt{ProgressPrinter}: per-step summary,
  \item \texttt{FiberRecorder}: fiber-level data,
  \item \texttt{ElementRecorder}: element-level data,
  \item \texttt{ResponseRecorder}: time-history at selected nodes,
  \item \texttt{VTMExporter}: VTK output for visualization,
  \item \texttt{DamageTracker}: damage ranking (every 10 steps),
  \item \texttt{FiberHysteresisRecorder}: hysteresis curves export.
\end{enumerate}
%
After the analysis, the \texttt{DamageTracker} prints the top-10
most damaged elements, and the \texttt{FiberHysteresisRecorder}
exports hysteresis CSV files for the extreme concrete and steel
fibers.

\subsection{Time Integration}

The implicit Generalized-$\alpha$ method
(\texttt{TSALPHA2} in PETSc) is used with:
%
\begin{itemize}
  \item $\Delta t = 0.005$~s,
  \item Rayleigh damping ($\alpha_M$, $\beta_K$ targeting
    $\xi = 5\%$ for the first two modes),
  \item total duration $T = 10.0$~s (matching the record length).
\end{itemize}


% ══════════════════════════════════════════════════════════════════════════
\section{Test Suite}
\label{sec:seismic-tests}

\texttt{tests/test\_seismic\_infrastructure.cpp} validates the full
infrastructure stack with nine test functions and 33~individual
assertions:

\begin{center}
\begin{tabular}{clc}
  \toprule
  \# & Test & Assertions \\
  \midrule
  1 & Parse El Centro record & 5 \\
  2 & \texttt{TimeFunction} evaluation & 5 \\
  3 & Scale factor ($g \to$ m/s$^2$) & 1 \\
  4 & \texttt{section\_snapshots()} beam fibers & 4 \\
  5 & \texttt{section\_snapshots()} elastic beam & 3 \\
  6 & \texttt{MaxStrainDamageCriterion} & 3 \\
  7 & \texttt{CallableDamageCriterion} & 5 \\
  8 & \texttt{GroundMotionBC} integration & 5 \\
  9 & \texttt{section\_snapshots()} for shells & 2 \\
  \midrule
  & Total & 33 \\
  \bottomrule
\end{tabular}
\end{center}

Tests~4--6 use a standalone \texttt{BeamFixture3D} that
constructs beam elements without a \texttt{Domain}:
%
\begin{lstlisting}[language=C++,basicstyle=\ttfamily\small]
struct BeamFixture3D {
    Node<3> n0{0, 0.0, 0.0, 0.0};
    Node<3> n1{1, 0.0, 0.0, 3.0};
    LagrangeElement3D<2> element{...};
    GaussLegendreCellIntegrator<2> integrator{};
    ElementGeometry<3> geom{element, integrator};
    BeamFixture3D() { geom.set_sieve_id(0); }
};
\end{lstlisting}
%
Test~9 uses an analogous \texttt{ShellFixture3D} with
four nodes and a \texttt{LagrangeElement3D<2,2>}.


% ══════════════════════════════════════════════════════════════════════════
\section{Build System Integration}
\label{sec:seismic-build}

\subsection{CMake Targets}

Two targets were added to \texttt{CMakeLists.txt}:
%
\begin{itemize}
  \item \texttt{fall\_n\_seismic}: the five-story RC building
        example (\texttt{main\_seismic.cpp}), linked against
        MPI, PETSc, Eigen3, VTK, OpenMP, and matplot++.
  \item \texttt{fall\_n\_seismic\_infra\_test}: the unit test
        executable, registered with CTest (label: \texttt{unit}).
\end{itemize}

\subsection{Header Organization}

\begin{itemize}
  \item \texttt{src/model/GroundMotionRecord.hh} is included
        by \texttt{header\_files.hh}.
  \item \texttt{src/analysis/DamageCriterion.hh} and
        \texttt{src/analysis/FiberHysteresisRecorder.hh}
        are included via \texttt{src/analysis/Analysis.hh}.
\end{itemize}


% ══════════════════════════════════════════════════════════════════════════
\section{Design Rationale}
\label{sec:seismic-rationale}

\subsection{Why a Virtual on StructuralElement?}

The alternative to adding \texttt{section\_snapshots()} would be
runtime type-switching with \texttt{as<T>()} for every concrete
element type.  This violates the Open--Closed Principle: every new
element type would require updating every observer.  The virtual
approach confines knowledge of ``how to extract sections'' to
the element's \texttt{Model<T>} specialization.

The \texttt{if constexpr (requires \{ element\_.sections(); \})}
guard uses C++20 concepts at compile time to enable the override
only for elements that have sections, making the virtual a
zero-overhead pass-through for element types that do not.

\subsection{Why Strategy Pattern for Damage?}

Seismic damage metrics vary widely:
\begin{itemize}
  \item maximum strain (implemented as default),
  \item Park--Ang index~\cite{parkang1985},
  \item inter-story drift ratio,
  \item energy-based criteria,
  \item user-defined combinations.
\end{itemize}
Making the criterion a pluggable strategy allows researchers to
inject domain-specific metrics \emph{without modifying library code}.
The \texttt{CallableDamageCriterion} enables one-line injection
via lambda.

\subsection{Why a Separate Hysteresis Recorder?}

Hysteresis curves require per-fiber, per-step time-series
$(\varepsilon(t), \sigma(t))$.  This is fundamentally different
from the snapshot-based \texttt{FiberRecorder} (latest step only)
or the scalar \texttt{DamageTracker} (peak ranking only).
A dedicated observer avoids polluting existing components
with time-history storage responsibilities.


% ══════════════════════════════════════════════════════════════════════════
\section{Summary of Files}
\label{sec:seismic-file-summary}

\begin{center}
\begin{tabular}{lp{8.5cm}}
  \toprule
  File & Purpose \\
  \midrule
  \texttt{data/input/earthquakes/el\_centro\_1940\_ns.dat}
    & El~Centro 1940 NS record (406 points) \\
  \texttt{src/model/GroundMotionRecord.hh}
    & Parser, accessors, \texttt{TimeFunction} conversion \\
  \texttt{src/analysis/DamageCriterion.hh}
    & \texttt{DamageCriterion} ABC, \texttt{MaxStrain},
      \texttt{Callable}, \texttt{DamageTracker} \\
  \texttt{src/analysis/FiberHysteresisRecorder.hh}
    & Fiber hysteresis tracking and CSV export \\
  \texttt{src/elements/StructuralElement.hh}
    & Added \texttt{section\_snapshots()} virtual \\
  \texttt{src/analysis/Observers.hh}
    & Wired \texttt{FiberRecorder}/\texttt{ElementRecorder}
      via \texttt{section\_snapshots()} \\
  \texttt{main\_seismic.cpp}
    & Five-story corotational RC building example \\
  \texttt{tests/test\_seismic\_infrastructure.cpp}
    & 9 tests, 33 assertions \\
  \bottomrule
\end{tabular}
\end{center}
